% Generated draft: requires mathematical and source review.
\documentclass[11pt]{article}
\usepackage[margin=25mm]{geometry}
\usepackage{fontspec}
\setmainfont{Latin Modern Roman}
\usepackage{amsmath,amssymb,mathtools,unicode-math}
\setmathfont{Latin Modern Math}
\usepackage{enumitem,needspace,xurl,hyperref}
\setlength{\parindent}{0pt}
\setlength{\parskip}{5pt}
\setlength{\emergencystretch}{3em}
\newcommand{\N}{\mathbb{N}}
\newcommand{\Z}{\mathbb{Z}}
\newcommand{\Q}{\mathbb{Q}}
\newcommand{\R}{\mathbb{R}}
\newcommand{\C}{\mathbb{C}}
\newcommand{\F}{\mathbb{F}}
\newcommand{\A}{\mathbb{A}}
\newcommand{\PP}{\mathbb{P}}
\newcommand{\E}{\mathbb{E}}
\newcommand{\eps}{\varepsilon}
\newcommand{\dd}{\,\mathrm{d}}
\begin{document}
\section*{UnsolvedMath ID: 30006536}
\textbf{Source problem code:} catalogue entry
\subsection{Definitions and mathematical statement}
The supplied record for UnsolvedMath permanent ID \(30006536\), source code \texttt{catalogue entry}, is titled ``Combinatorics Contents Fragment.'' Its statement consists of a table-of-contents fragment listing unrelated talks, preceded by an incomplete entry ending with ``problem'' and the page number \(15\).

No mathematical problem can be stated precisely from this fragment. It specifies no target proposition, mathematical objects, hypotheses, domains, boundary conditions, or quantifiers. The listed talk titles do not supply definitions or conjecture statements, and the record does not identify any one talk as the intended problem. Consequently, no mathematical definitions or assertion are reconstructed here.

The supplied \texttt{statement\_status} is \texttt{unrecoverable}. The supplied statement-verification note reports that the official report confirms the extraction is a table-of-contents fragment spanning unrelated talks and states no problem; it also says that no separate clean statement is promoted and that the original record remains publicly visible for provenance. These are supplied curation claims, not independently verified findings of this draft. The registry labels the record ``open,'' but its supplied literature assessment explicitly says that no specific proposition can be assessed as solved or disproved.
\subsection{Short English statement}
This record contains contents entries for several talks, without an identifiable mathematical problem. A precise statement requires recovery of the intended passage and its hypotheses. The supplied ``open'' label does not establish the status of any particular proposition.
\subsection{Sources}
\begin{itemize}
\item Local UnsolvedMath registry: permanent ID \(30006536\), exact source code \texttt{catalogue entry}, title ``Combinatorics Contents Fragment,'' collection ``Oberwolfach Reports Open Problems.'' No problem proposer is identified in the supplied record.
\item Supplied report citation: \emph{Combinatorics} (2026), Oberwolfach Reports, DOI 10.4171/owr/2026/1, pp.~5--79. Supplied source URL and literature-source link, labeled ``Oberwolfach Report 2026/1'': \url{https://doi.org/10.4171/owr/2026/1}. The background describes Oberwolfach Reports as documenting mathematical workshops and their open-problem discussions.
\item The supplied fragment begins with an incomplete contents entry ending in ``problem,'' p.~15; its title and attribution cannot be recovered from the record.
\item Michael Krivelevich, joint with Sahar Diskin, Itay Markbreit, and Maksim Zhukovskii: ``A large hole in pseudo-random graphs,'' listed on p.~17.
\item Matija Bucić, joint with K. Hendrey, B. Mohar, R. Steiner, and L. Yepremyan: ``Long cycles in vertex transitive digraphs,'' listed on p.~18.
\item Matthew Kwan, joint with Zach Hunter and Lisa Sauermann: ``Exponential anticoncentration of the permanent,'' listed on p.~20.
\item Benny Sudakov, joint with Zach Hunter, Aleksa Milojević, and István Tomon: ``\(C_4\)-free subgraphs of high degree with geometric applications,'' listed on p.~23.
\item Anubhab Ghosal, joint with Peter Keevash: ``On the largest \(k\)-product-free subsets of the Symmetric and Alternating Groups,'' listed on p.~27.
\item Noam Lifshitz, joint with Shai Evra and Guy Kindler: ``Polynomial Bogolyubov for Special Linear Groups via Tensor Rank,'' listed on p.~29.
\item Peter Allen, joint with Julia Böttcher, Jozef Skokan, and Benny Sudakov: ``A bipartite Bollobás-Eldridge-Catlin conjecture,'' listed on p.~31.
\item Matthew Jenssen, joint with Marcelo Campos, Marcus Michelen, Florian Pfender, and Julian Sahasrabudhe: ``A polynomial improvement for the odd cycle--complete Ramsey numbers,'' listed on p.~32.
\item Oliver Janzer, joint with Maria Axenovich, Wouter Cames van Batenburg, Lukas Michel, and Mathieu Rundström: ``An improved upper bound for the multicolour Ramsey number of odd cycles,'' listed on p.~34.
\item Jeff Kahn, joint with Max Aires: ``More on balancing extensions,'' listed on p.~36. The preceding names and titles are contents-entry attributions, not identified proposers or statements of this problem.
\item Supplied literature assessment, checked on \texttt{2026-08-22}: status ``open''; the exact record is a table-of-contents fragment rather than a mathematical problem statement, so no specific proposition can be assessed as solved or disproved. The supplied background lists the report above under ``Sources checked.'' This assessment is supplied dated curation derived from the supplied OWR curation, not independently verified here; time-sensitive status and cited claims require independent verification.
\end{itemize}
\end{document}
